\documentclass{article}
\usepackage{amsmath}
\usepackage{amssymb}
\usepackage{array}

\begin{document}
	
	\section*{Christoffel Symbols: Unified View (Polar Coordinates)}
	
	\subsection*{1. Coordinate System}
	
	\[
	(r,\theta)
	\]
	
	Metric:
	\[
	g_{ij} =
	\begin{pmatrix}
		1 & 0 \\
		0 & r^2
	\end{pmatrix}
	\quad
	g^{ij} =
	\begin{pmatrix}
		1 & 0 \\
		0 & \frac{1}{r^2}
	\end{pmatrix}
	\]
	
	\subsection*{2. Summary of Nonzero Christoffel Symbols}
	
	\[
	\Gamma^r_{\theta\theta} = -r
	\quad
	\Gamma^\theta_{r\theta} = \Gamma^\theta_{\theta r} = \frac{1}{r}
	\]
	
	\subsection*{3. Comparison of the Three Approaches}
	
	\renewcommand{\arraystretch}{1.5}
	
	\begin{tabular}{| m{3.8cm} | m{5.5cm} | m{5.5cm} |}
		\hline
		\textbf{Method} & \textbf{Key Formula} & \textbf{Interpretation} \\
		\hline
		
		Metric (algebraic) &
		\[
		\Gamma^k_{ij} =
		\frac{1}{2} g^{kl}
		(\partial_i g_{jl} + \partial_j g_{il} - \partial_l g_{ij})
		\]
		&
		Derived purely from the metric. Measures how the inner product changes across space. \\
		
		\hline
		
		Basis derivatives (geometric) &
		\[
		\frac{\partial \mathbf{e}_i}{\partial x^j}
		=
		\Gamma^k_{ij} \mathbf{e}_k
		\]
		&
		Christoffel symbols describe how the coordinate basis vectors change from point to point. \\
		
		\hline
		
		Covariant derivative / projection &
		\[
		\Gamma^k_{ij}
		=
		\mathbf{e}^k \cdot \frac{\partial \mathbf{e}_i}{\partial x^j}
		\]
		&
		Christoffel symbols are components of the derivative of a basis vector projected onto the dual basis. \\
		
		\hline
		
	\end{tabular}
	
	\subsection*{4. Equivalence of the Methods}
	
	All three approaches give identical results because:
	
	\[
	g_{ij} = \mathbf{e}_i \cdot \mathbf{e}_j
	\]
	
	Taking derivatives:
	
	\[
	\partial_k g_{ij}
	=
	\left(\frac{\partial \mathbf{e}_i}{\partial x^k}\right)\cdot \mathbf{e}_j
	+
	\mathbf{e}_i \cdot \left(\frac{\partial \mathbf{e}_j}{\partial x^k}\right)
	\]
	
	Substituting:
	\[
	\frac{\partial \mathbf{e}_i}{\partial x^k} = \Gamma^m_{ik}\mathbf{e}_m
	\]
	
	yields the metric formula.
	
	\subsection*{5. Geodesic Equations}
	
	A geodesic is defined by:
	\[
	\frac{d^2 x^k}{d\lambda^2}
	+
	\Gamma^k_{ij}
	\frac{dx^i}{d\lambda}
	\frac{dx^j}{d\lambda}
	=0
	\]
	
	\subsection*{6. Explicit Equations in Polar Coordinates}
	
	\subsubsection*{Radial equation ($k=r$)}
	
	\[
	\frac{d^2 r}{d\lambda^2}
	+
	\Gamma^r_{\theta\theta}
	\left(\frac{d\theta}{d\lambda}\right)^2
	=0
	\]
	
	\[
	\frac{d^2 r}{d\lambda^2}
	- r \left(\frac{d\theta}{d\lambda}\right)^2
	=0
	\]
	
	\subsubsection*{Angular equation ($k=\theta$)}
	
	\[
	\frac{d^2 \theta}{d\lambda^2}
	+
	2 \Gamma^\theta_{r\theta}
	\frac{dr}{d\lambda}
	\frac{d\theta}{d\lambda}
	=0
	\]
	
	\[
	\frac{d^2 \theta}{d\lambda^2}
	+
	\frac{2}{r}
	\frac{dr}{d\lambda}
	\frac{d\theta}{d\lambda}
	=0
	\]
	
	\subsection*{7. Interpretation}
	
	\begin{itemize}
		\item The radial equation:
		\[
		\frac{d^2 r}{d\lambda^2} - r\dot{\theta}^2 = 0
		\]
		represents the familiar ``centrifugal'' term.
		
		\item The angular equation:
		\[
		\frac{d^2 \theta}{d\lambda^2}
		+ \frac{2}{r}\dot{r}\dot{\theta} = 0
		\]
		encodes conservation of angular momentum.
		
	\end{itemize}
	
	\subsection*{8. Final Insight}
	
	Geodesics are curves whose tangent vector is parallely transported:
	\[
	\nabla_{\dot{\gamma}} \dot{\gamma} = 0
	\]
	
	Christoffel symbols act as the correction terms that compensate for the curvature of the coordinate system.
	
\end{document}
